%%%%%%%%%%%%%%%%%%%%%%%%%%%%%%%%%%%%%%%%%%%%%%%%%%%%%%%%%%%%%%%%%%%%%%%%%%
%%%%%%%%%%%%%%%%%%%%%%%%%%%       Aro       %%%%%%%%%%%%%%%%%%%%%%%%%%%%%%%
%%%%%%%%%%%%%%%%%%%%%%%%%%%%%%%%%%%%%%%%%%%%%%%%%%%%%%%%%%%%%%%%%%%%%%%%%%

% Dois círculos de centro na origem, de raios r < R, num plano coordenado,
% com parte da região entre eles (coroa circular) sombreada.
% Uso: \Aro{eixo horizontal}{eixo vertical}{r}{R}{q}
% Ex.: \Aro{x}{y}{1}{2}{1}  % parte do aro no 1º quadrante do plano xy
%      \Aro{y}{z}{1}{3}{t}  % metade de cima do aro no plano yz
% q = 1, 2, 3, 4 : um quadrante
% q = t, b, l, r : metade de cima (top), de baixo (bottom), esquerda (left) ou direita (right)
% q = o          : aro todo

% Comando auxiliar, com a região definida explicitamente
% #1, #2: eixos (horizontal e vertical)   #3, #4: raios menor (r) e maior (R)
% #5, #6: ângulos inicial (i) e final (j)
\newcommand{\AroAux}[6]{%
\begin{tikzpicture}[scale=1]
  \begin{axis}[
      xlabel={$#1$}, ylabel={$#2$},
      axis equal image,                % mantém escala para círculos não ficarem ovalados
      xtick={#3,#4}, ytick=\empty,
      x tick label style={anchor=south west}, % rótulos acima e à direita do eixo horizontal
      every axis x label/.style={at={(ticklabel* cs:0.99)}, anchor=west},
      every axis y label/.style={at={(ticklabel* cs:0.99)}, anchor=south},
    ]

    % --- Curvas ---
    \addplot[thick, blue, domain=0:360, samples=360]({#3*cos(x)},{#3*sin(x)});
    \addplot[thick, blue, domain=0:360, samples=360]({#4*cos(x)},{#4*sin(x)});

    % --- Região sombreada ---
    \path[fill=gray!30, opacity=0.6]
      (axis cs:{#4*cos(#5)},{#4*sin(#5)})              % Ponto inicial, no círculo maior
      arc[start angle=#5, end angle=#6, radius=#4]      % Arco do círculo maior, de i a j
      -- (axis cs:{#3*cos(#6)},{#3*sin(#6)})           % Segmento até o círculo menor
      arc[start angle=#6, end angle=#5, radius=#3]      % Arco do círculo menor, de volta de j a i
      -- cycle;                                         % Fecha no ponto inicial

  \end{axis}
\end{tikzpicture}%
}

% Tabela de q -> (i, j). Cada entrada recebe #1, #2 (eixos), #3 (r) e #4 (R).
\makeatletter
\@namedef{Aro@1}#1#2#3#4{\AroAux{#1}{#2}{#3}{#4}{0}{90}}
\@namedef{Aro@2}#1#2#3#4{\AroAux{#1}{#2}{#3}{#4}{90}{180}}
\@namedef{Aro@3}#1#2#3#4{\AroAux{#1}{#2}{#3}{#4}{180}{270}}
\@namedef{Aro@4}#1#2#3#4{\AroAux{#1}{#2}{#3}{#4}{270}{360}}
\@namedef{Aro@t}#1#2#3#4{\AroAux{#1}{#2}{#3}{#4}{0}{180}}
\@namedef{Aro@l}#1#2#3#4{\AroAux{#1}{#2}{#3}{#4}{90}{270}}
\@namedef{Aro@b}#1#2#3#4{\AroAux{#1}{#2}{#3}{#4}{180}{360}}
\@namedef{Aro@r}#1#2#3#4{\AroAux{#1}{#2}{#3}{#4}{-90}{90}}
\@namedef{Aro@o}#1#2#3#4{\AroAux{#1}{#2}{#3}{#4}{0}{360}}

\newcommand{\Aro}[5]{%
  \@ifundefined{Aro@#5}%
    {\PackageError{Aro}{q = '#5' invalido}{Use 1, 2, 3, 4, t, b, l, r ou o.}}%
    {\@nameuse{Aro@#5}{#1}{#2}{#3}{#4}}%
}
\makeatother
